% !TeX root = ../papers/nc-p3.tex
\section{四行列恒等式の全単純の均衡性と半素商のNoether性}
\label{sec:four-matrix-balanced-simples}
\begin{theorem}\label{thm:four-matrix-balanced-simples}
全四隅が四行列の恒等式を満たす任意の有限四辺の標準正則窓で、
全非頂点単純の次元ベクトルは $\delta,2\delta,3\delta,4\delta$ のいずれかである。
さらに根基を割った代数は左右Noetherである。一隅半素なら全窓が中心有限かつ合成整域となる。
これは十五次元閉包の外も含む。
\end{theorem}
\begin{proof}[最後の混合次元の排除]
一つの頂点が4に達すれば四行列の中心多項式と隅中心の延長で均衡となる。
薄いもの以外は全頂点次元が2以上である。前節から残る
$S$ の次元 $(2,2,2,3)$ を仮定する。別の一般次数倍率 $T$ を取る。
両者の第三頂点次元はともに3なので両向きの第一Extは零、第二ExtはEuler式で次元2である。
頂点への向きを前節の $a_i,b_i$ とすれば $h=(-4,6,-4,2)$。
零番で $(a_0,b_0)\le(4,1)$、二番で $(a_2,b_2)\le(1,4)$ であり、各和は4以上必要である。
$a_0b_0+a_2b_2\le2$ なので $(a_0,b_0)=(4,0),(a_2,b_2)=(0,4)$ しかない。
これらは $T$ についても同じで、各二位置で頂点との第二Extは両向きに零である。
一番では $a_1,b_1\le1$。

四つの独立な拡張を全て使い
\[
 0\to k_0^4\to U\to S\to0,\qquad
 0\to T\to V\to k_2^4\to0
\]
を取る。$U$ の頭は $S$ 一つ、$V$ のsocleは $T$ 一つである。例えば $U$ から
$k_0$ への射は四拡張の関係を要求して零となる。
頭 $S$ だけを持つ加群 $P$ とsocle $T$ だけを持つ加群 $Q$ に対し、
本証明で使うものは $P$ 内に $T$、$Q$ 内に $S$ を含まず、$S,T$ は各一回である。
この場合 $\operatorname{Ext}^1(P,Q)=0$ である。非分裂拡張の有限像で持ち上げた
$S$ の行列ブロックが生成する部分加群は $P$ 全体へ全射し、核が零なら分裂する。
非零なら $Q$ 内の $T$ を含む。従って $S$ のブロックから $T$ のブロックへの
非零矩形成分が存在するが、第三頂点だけで $3+3>4$ となり不可能である。
同じ理由で補題\ref{lem:pi-yoneda-bridges}の中間 $k_1$ による積の単射性は、
端点を $(U,T),(S,V),(U,V)$ に置き換えても成立する。
実際中間の重複 $k_1^m$ を $P$ の上に付けた拡張で用いた $m$ 類が独立なら、
その加群に $k_1$ の商はなく、頭は依然 $S$ 一つである。下半分もsocleは $T$ 一つである。
積が零なら組める加群で、頭から中間全体へ到達し、中間の非零成分からsocleへ到達する。
中間 $k_1$ はどちらの端の組成因子でもないのでその持上げブロックを挿入でき、
二端を結ぶ非零矩形を得る。この議論は $k_1$ 自身が第三頂点で零であっても、
できた合成の両端だけを第三頂点に制限するので成立する。

完全列より $\dim\operatorname{Ext}^2(U,T),\dim\operatorname{Ext}^2(S,V)\le2$。
一方、$\dim\operatorname{Ext}^1(k_0,k_1)=\dim\operatorname{Ext}^1(k_1,k_2)=4$、
$\dim\operatorname{Ext}^2(S,k_1)=6+a_1+b_1$ だから
\[
   \dim\operatorname{Ext}^1(U,k_1)\ge10-b_1,\qquad
   \dim\operatorname{Ext}^1(k_1,V)\ge10-a_1 .
\]
例えば第一は $a_1+16-(6+a_1+b_1)$ である。積の単射性をまず $(U,T),(S,V)$ で使うと
$(10-b_1)b_1\le2$, $a_1(10-a_1)\le2$ だから $a_1=b_1=0$。
次に $(U,V)$ で使うと第二Extの次元は少なくとも100になる。

しかし $U,V$ の組成因子は交わらない。両向きHomと第四Extが零で、上で第一Extの片向きも除いた。
逆向き第一Extは組成列で上から評価すると高々32である。
寄与し得るのは $T$ から $k_0^4$ の16と $k_2^4$ から $S$ の16のみで、
$k_2,k_0$ は隣接せず、$T,S$ の第一Extも零である。Euler数は
\[
          \chi(U,V)=2-4\cdot4-4\cdot4+16\cdot6=66 .
\]
CY双対により第二Extは高々 $66+32=98$ となり矛盾する。
\end{proof}

\begin{proof}[全半素商の有限性]
斉次素商のtrace環の一般単純を特殊化し、半単純型が頂点と非頂点を混ぜると仮定する。
頂点次元は高々4で、非頂点因子は今や均衡している。従ってその中で非薄なものは重複込みで高々一つである。
二つあれば $2\delta+2\delta$ だけで頂点の入る余地はない。
一つなら薄い因子はいくつあっても頂点との全第一Extが零であり、残る $2\delta$ または $3\delta$ の一つも
定理\ref{thm:one-sided-balanced-small-simples}により頂点との拡張は高々一方向である。
従って頂点群と非頂点群の間の片側三角化により、前節の証明と同じく一般絶対単純のこの型への特殊化は不可能。
非原点のtrace指標で頂点単純が出ないことが分かる。
よって各斉次素商は自身のNoetherなtrace環に有限中心化拡大を持ち、左右Noetherである。
根基は斉次で、半素商は有限個の斉次極小素商の積の部分加群として左右Noetherになる。
一隅半素の場合は前節と同じく折り畳み全体が半素なので、既出のNoether正則PI窓の定理が結論を与える。
\end{proof}
